\section{Continuous histories, dropout and separately learned policies}
\label{s5:history-censoring}
Finite tables are useful only when their state captures the information
that jointly determines action and outcome. A servicing system may instead
observe utilization, recent payment movements, earlier assistance and the
number of previous interventions. Treat the resulting history $H_t$ as
continuous. Let $D_t$ indicate that reward $R_t$ and next history $H_{t+1}$
are observed, conditional on remaining observed before decision $t$.
Censoring is monotone; default itself is an outcome, not dropout. The
target intervention replaces action probabilities by $\pi_t(a\mid h)$
and ensures follow-up. This joint intervention requires an argument that
follow-up does not itself change the customer process, or a model for how
it does. It is not the mean among customers who happened to remain visible.

Write $e_t(a\mid h)$ for the historical action probability and
$c_t(h,a)$ for conditional survival of observation. Consistency,
sequential exchangeability for actions and censoring, and positive
$e_tc_t$ on the policy's support identify
\begin{align}
 Q_t(h,a)&=\E[R_t+\gamma V_{t+1}(H_{t+1})\mid H_t=h,A_t=a,D_t=1],\\
 V_t(h)&=\sum_a\pi_t(a\mid h)Q_t(h,a),\qquad V_H=0,\qquad
 v_\pi=\E[V_0(H_0)].
 \label{s5:censor-recursion}
\end{align}
Exchangeability must concern the future histories used by later decisions,
as well as the terminal reward. Under the observed-data model the trajectory
score is
\begin{equation}
 U_\pi=V_0(H_0)+\sum_{t=0}^{H-1}\gamma^t W_t
       \{R_t+\gamma V_{t+1}(H_{t+1})-Q_t(H_t,A_t)\},\quad
 W_t=\prod_{s=0}^t\frac{\pi_s(A_s\mid H_s)D_s}
 {e_s(A_s\mid H_s)c_s(H_s,A_s)}.
 \label{s5:censor-score}
\end{equation}
Every summand after dropout is defined to be zero without accessing the
unobserved reward/history. The implemented check changes those hidden
values arbitrarily and verifies an identical score. This is a substantive
observation constraint, not merely a missing-value imputation convention.

\subsection{Telescoping, the remainder and the machine-learning obligation}
Conditional on the fitted nuisances, put
$\hat r_t=e_tc_t/(\hat e_t\hat c_t)$ and
$\delta_t=\hat Q_t-Q_t$, with
$\delta V_t=\sum_a\pi_t(a\mid H_t)\delta_t(H_t,a)$.
Changing measure one step at a time in the fitted score gives
\begin{align*}
 \E[\hat U_\pi]-v_\pi
 &=\E[\delta V_0]
   +\sum_t\gamma^t\E^{\pi}
       \left[\prod_{s\le t}\hat r_s
       \{\gamma\delta V_{t+1}-\delta_t\}\right]\\
 &=\sum_t\gamma^t\E^{\pi}
       \left[\prod_{s<t}\hat r_s(1-\hat r_t)\delta_t\right].
\end{align*}
Here $\E^\pi$ is the full-follow-up interventional path expectation; the
empty product is one. The second equality pairs the time-$t$ regression
error with the continuation error at time $t-1$, integrating the current
action under $\pi$. It exposes exactly why both the treatment and censoring
mechanisms belong in the weight. It also proves zero bias if every $Q_t$
is correct, or every product $e_tc_t$ is correct. More selective robustness
requires the appropriate backward consistency of the estimated continuation
regressions; it is not obtained by declaring an arbitrary collection of
models correct after fitting.

For fixed horizon, bounded policy ratios, and nuisance probabilities
bounded away from zero, Cauchy--Schwarz bounds this remainder by
\[
 C\sum_t\|\hat Q_t-Q_t\|_2
   \{\|\hat e_t-e_t\|_2+\|\hat c_t-c_t\|_2\}.
\]
The norms must dominate the relevant policy occupations; fixed overlap
provides the required change-of-measure bounds. Cross-fitting whole
customers controls the evaluation empirical process if the fitted score
converges in $L_2$. Product error $o_p(N^{-1/2})$, finite score variance and
an independent-trajectory CLT then give
$\sqrt N(\hat v_\pi-v_\pi)=N^{-1/2}\sum_i(U_{\pi,i}-v_\pi)+o_p(1)$.
For two policies use $U_{\pi_1,i}-U_{\pi_0,i}$ within each customer and its
variance. Summing two separately estimated variances would lose their
covariance. Fixed clipping is justified only if it is asymptotically inactive
under the asserted overlap; otherwise its bias or changed target must be
analyzed. A fitted probability with a good calibration plot does not prove
exchangeability, overlap in unobserved histories, or the product-rate bound.

The estimator is compatible with flexible regression, but its theorem is
conditional on those rates. Histogram boosting in the experiment is an
implemented nuisance choice, not a general proof of its rates for arbitrary
customer histories. Recent functional-history work makes a similar
distinction: a learned representation must preserve the outcome and mechanism
information needed by the score \citep[Definition 1, Assumptions 5--9,
Theorems 3--4]{ran2026}. Its approximation errors enter products alongside
estimation error. A representation with small prediction loss can still
discard a confounder. Its target concerns decisions based on observed
irregular measurements; a policy based on an unobserved complete trajectory
needs a further measurement model. The source's mean-aggregation Wald result
also does not automatically extend to its separate robust point aggregators.

\subsection{Selection, censoring and the fitted comparison}
An independent randomized pilot of 2,500 synthetic customers learns a
four-coefficient immediate-reward contrast. The resulting stochastic policy
uses action probability .7 where that fitted contrast exceeds .5 and .2
elsewhere; the reference policy always uses .2. This is a fixed learned
candidate after the pilot, not a proof of dynamic optimality. Conditional
on the pilot, the target is its actual policy value. An unconditional claim
about a policy-learning algorithm would require a different sampling
experiment that repeats the pilot as well. Regret relative to an optimal
policy is another estimand.

Evaluation uses 2,400 fresh customers, four decisions and discount .97.
Histories include a baseline covariate, continuous current and lagged risk,
previous action and cumulative action count. These summaries are sufficient
by the declared generator; no generic learned-state sufficiency is assumed.
The action and observation probabilities are fitted with logistic regressions;
arm-specific histogram boosting estimates continuation rewards backward.
Two whole-trajectory folds keep each evaluated customer's nuisance fits
outside that fold. A separate million-trajectory intervention simulation
gives contrast $3.5813$ with Monte Carlo standard error $.0040$, much smaller
than the evaluation interval standard error of about $.295$.

Across 120 repetitions under coarsening at random, the full score has bias
$-.0012$, RMSE $.2873$ and 95.8\% coverage at nominal 95\%. Ignoring
censoring gives bias $-.0199$ and 83.3\% coverage: the continuation model
is useful, yet its score variance is no longer an adequate inferential
construction. Direct recursion has mean bias $-.0376$; it is a point
challenger without a claimed interval. In the adverse generator, observation
also depends on the otherwise unseen reward innovation. The full fitted
score now has bias $-.1434$ and 89.2\% coverage; ignoring censoring covers
57.5\%. This is an observed failure of the identifying censoring condition,
not proof that a different learner will repair MNAR data. At 120 repetitions
the Monte Carlo standard error of coverage near 95\% is about .020.

Irregular opportunities require a further distinction. One may model the
timestamps and elapsed time as part of each observed history and retain the
natural visit clock; changing when visits occur is a different intervention.
The event-intensity likelihood and weights in Chapter~\ref{r3:continuous}
provide that second route. Interpolating missing months and treating them as
observed decisions supplies neither exchangeability nor the visit likelihood.
Thus the proposed servicing workflow first selects among observed-history
adjustment, an explicit controlled measurement model, and a proxy bridge.
Their identifying information differs even when all three output a policy
value in the same currency.
